\bibitem{budak2022sizing}
Budak AF, Gandara M, Shi W, Pan DZ, Sun N, Liu B. An Efficient Analog Circuit Sizing Method Based on Machine Learning Assisted Global Optimization. IEEE Trans. Comput.-Aided Des. Integr. Circuits Syst. 2022;41:1209-1221.
\href{https://doi.org/10.1109/tcad.2021.3081405}{https://doi.org/10.1109/tcad.2021.3081405}

\bibitem{rashid2024global}
Rashid R, Krishna K, George CP, Nambath N. Machine learning driven global optimisation framework for analog circuit design. Microelectronics Journal. 2024;151:106362.
\href{https://doi.org/10.1016/j.mejo.2024.106362}{https://doi.org/10.1016/j.mejo.2024.106362}

\bibitem{yin2023surrogate}
Yin S, Wang R, Zhang J, Liu X, Wang Y. Fast Surrogate-Assisted Constrained Multiobjective Optimization for Analog Circuit Sizing via Self-Adaptive Incremental Learning. IEEE Trans. Comput.-Aided Des. Integr. Circuits Syst. 2023;42:2080-2093.
\href{https://doi.org/10.1109/tcad.2022.3221694}{https://doi.org/10.1109/tcad.2022.3221694}

\bibitem{huang2021mleda}
Huang G, Hu J, He Y, Liu J, Ma M, Shen Z, et al. Machine Learning for Electronic Design Automation: A Survey. ACM Trans. Des. Autom. Electron. Syst. 2021;26:1-46.
\href{https://doi.org/10.1145/3451179}{https://doi.org/10.1145/3451179}

\bibitem{jang2023gnn}
Jang W, Myung S, Choe JM, Kim YG, Kim DS. TCAD Device Simulation With Graph Neural Network. IEEE Electron Device Lett. 2023;44:1368-1371.
\href{https://doi.org/10.1109/led.2023.3290930}{https://doi.org/10.1109/led.2023.3290930}

\bibitem{wang2023tcadcal}
Wang BC, Tang CX, Qiu MT, Chen W, Wang T, Xu JY, et al. A machine learning approach to TCAD model calibration for MOSFET. NUCL SCI TECH. 2023;34.
\href{https://doi.org/10.1007/s41365-023-01340-x}{https://doi.org/10.1007/s41365-023-01340-x}

\bibitem{li2025devmodel}
Li X, Wu Z, Rzepa G, Karner M, Xu H, Wu Z, et al. Overview of emerging semiconductor device model methodologies: From device physics to machine learning engines. Fundamental Research. 2025;5:2149-2160.
\href{https://doi.org/10.1016/j.fmre.2024.01.010}{https://doi.org/10.1016/j.fmre.2024.01.010}

\bibitem{liu2025sann}
Liu Y, Liu X, Zhou Y, Cao P. S-ANN: Synchronous TCAD device simulation of FinFET using Artificial Neural Network. Microelectronics Journal. 2025;159:106630.
\href{https://doi.org/10.1016/j.mejo.2025.106630}{https://doi.org/10.1016/j.mejo.2025.106630}

\bibitem{maniyar2022sidewall}
Maniyar A, Srinivas PSTN, Tiwari PK, Chang-Liao KS. Impact of Process-Induced Inclined Sidewalls on Gate-Induced Drain Leakage (GIDL) Current of Nanowire GAA MOSFETs. IEEE Trans. Electron Devices. 2022;69:4815-4820.
\href{https://doi.org/10.1109/ted.2022.3194109}{https://doi.org/10.1109/ted.2022.3194109}

\bibitem{lee2022stress}
Lee K, Kaczer B, Kruv A, Gonzalez M, Eneman G, Okudur OO, et al. Gate-Induced-Drain-Leakage (GIDL) in CMOS Enhanced by Mechanical Stress. IEEE Trans. Electron Devices. 2022;69:2214-2217.
\href{https://doi.org/10.1109/ted.2022.3154341}{https://doi.org/10.1109/ted.2022.3154341}

\bibitem{thakur2023gidl}
Thakur A, Dhiman R. Comprehensive study of gate induced drain leakage in nanowire and nanotube junctionless FETs using Si1-xGex source/drain. AEU - International Journal of Electronics and Communications. 2023;167:154668.
\href{https://doi.org/10.1016/j.aeue.2023.154668}{https://doi.org/10.1016/j.aeue.2023.154668}

\bibitem{kaushal2022ncfinfet}
Kaushal S, Rana AK. Modeling and analysis of gate-induced drain leakage current in negative capacitance junctionless FinFET. J Comput Electron. 2022;21:1229-1238.
\href{https://doi.org/10.1007/s10825-022-01930-9}{https://doi.org/10.1007/s10825-022-01930-9}

\bibitem{bashir2024btbt}
Bashir MY, Jaiswal AK, Patel SD, Sahay S. Revisiting Lateral-BTBT Gate-Induced Drain Leakage in Nanowire FETs for 1T-DRAM. IEEE Trans. Electron Devices. 2024;71:2714-2720.
\href{https://doi.org/10.1109/ted.2024.3365773}{https://doi.org/10.1109/ted.2024.3365773}

\bibitem{kwak2024gidl}
Kwak B, Lee K, Kim S, Kwon D. Suppression of Gate-Induced-Drain-Leakage Utilizing Local Polarization in Ferroelectric-Gate Field-Effect Transistors for DRAM Applications. IEEE Electron Device Lett. 2024;45:813-816.
\href{https://doi.org/10.1109/led.2024.3370592}{https://doi.org/10.1109/led.2024.3370592}

\bibitem{kim2021trap}
Kim KY, Min KK, Park BG. Trap-Induced Data-Retention-Time Degradation of DRAM and Improvement Using Dual Work-Function Metal Gate. IEEE Electron Device Lett. 2021;42:38-41.
\href{https://doi.org/10.1109/led.2020.3037640}{https://doi.org/10.1109/led.2020.3037640}

\bibitem{kim2021rdf}
Kim KY, Park BG. Effect of Random Dopant Fluctuation on Data Retention Time Distribution in DRAM. IEEE Trans. Electron Devices. 2021;68:5572-5577.
\href{https://doi.org/10.1109/ted.2021.3108743}{https://doi.org/10.1109/ted.2021.3108743}

\bibitem{bepary2022aging}
Bepary MK, Talukder BMSB, Rahman MT. DRAM Retention Behavior with Accelerated Aging in Commercial Chips. Applied Sciences. 2022;12:4332.
\href{https://doi.org/10.3390/app12094332}{https://doi.org/10.3390/app12094332}

\bibitem{park2025capleak}
Park DS, Lee EO, Lee JM, Cho E, Choi B. Suppression of Capacitor Leakage Through Thermal Budget Control in DRAM With ZrO 2 -Based Dielectrics. IEEE Access. 2025;13:35019-35028.
\href{https://doi.org/10.1109/access.2024.3502036}{https://doi.org/10.1109/access.2024.3502036}

\bibitem{du2025crdram}
Du H, Zhu H, Chen S, Kang Y. CR-DRAM: Improving DRAM Refresh Energy Efficiency With Inter-Subarray Charge Recycling. IEEE Trans. VLSI Syst. 2025;33:21-34.
\href{https://doi.org/10.1109/tvlsi.2024.3445631}{https://doi.org/10.1109/tvlsi.2024.3445631}

\bibitem{golman2021refresh}
Golman R, Nachum N, Cohen T, Giterman R, Teman A. Refresh Algorithm for Ensuring 100\% Memory Availability in Gain-Cell Embedded DRAM Macros. IEEE Access. 2021;9:105831-105840.
\href{https://doi.org/10.1109/access.2021.3099970}{https://doi.org/10.1109/access.2021.3099970}

\bibitem{nguyen2021zem}
Nguyen DT, Ho NM, Le MS, Wong WF, Chang IJ. ZEM: Zero-Cycle Bit-Masking Module for Deep Learning Refresh-Less DRAM. IEEE Access. 2021;9:93723-93733.
\href{https://doi.org/10.1109/access.2021.3088893}{https://doi.org/10.1109/access.2021.3088893}

\bibitem{gupta2021tfet}
Gupta N, Makosiej A, Shrimali H, Amara A, Vladimirescu A, Anghel C. Tunnel FET Negative-Differential-Resistance Based 1T1C Refresh-Free-DRAM, 2T1C SRAM and 3T1C CAM. IEEE Trans. Nanotechnology. 2021;20:270-277.
\href{https://doi.org/10.1109/tnano.2021.3061607}{https://doi.org/10.1109/tnano.2021.3061607}

\bibitem{zhang2022batchbo}
Zhang S, Yang F, Yan C, Zhou D, Zeng X. An Efficient Batch-Constrained Bayesian Optimization Approach for Analog Circuit Synthesis via Multiobjective Acquisition Ensemble. IEEE Trans. Comput.-Aided Des. Integr. Circuits Syst. 2022;41:1-14.
\href{https://doi.org/10.1109/tcad.2021.3054811}{https://doi.org/10.1109/tcad.2021.3054811}

\bibitem{touloupas2022locomobo}
Touloupas K, Sotiriadis PP. LoCoMOBO: A Local Constrained Multiobjective Bayesian Optimization for Analog Circuit Sizing. IEEE Trans. Comput.-Aided Des. Integr. Circuits Syst. 2022;41:2780-2793.
\href{https://doi.org/10.1109/tcad.2021.3121263}{https://doi.org/10.1109/tcad.2021.3121263}

\bibitem{chen2022highdimbo}
Chen C, Wang H, Song X, Liang F, Wu K, Tao T. High-Dimensional Bayesian Optimization for Analog Integrated Circuit Sizing Based on Dropout and gm/ID Methodology. IEEE Trans. Comput.-Aided Des. Integr. Circuits Syst. 2022;41:4808-4820.
\href{https://doi.org/10.1109/tcad.2022.3147431}{https://doi.org/10.1109/tcad.2022.3147431}

\bibitem{saglican2023moead}
Saglican E, Afacan E. MOEA/D vs. NSGA-II: A Comprehensive Comparison for Multi/Many Objective Analog/RF Circuit Optimization through a Generic Benchmark. ACM Trans. Des. Autom. Electron. Syst. 2023;29:1-23.
\href{https://doi.org/10.1145/3626096}{https://doi.org/10.1145/3626096}

\bibitem{abdollahi2024xai}
Abdollahi A, Li D, Deng J, Amini A. An explainable artificial-intelligence-aided safety factor prediction of road embankments. Engineering Applications of Artificial Intelligence. 2024;136:108854.
\href{https://doi.org/10.1016/j.engappai.2024.108854}{https://doi.org/10.1016/j.engappai.2024.108854}

\bibitem{homafar2022shap}
Homafar A, Nasiri H, Chelgani SC. Modeling coking coal indexes by SHAP-XGBoost: Explainable artificial intelligence method. Fuel Communications. 2022;13:100078.
\href{https://doi.org/10.1016/j.jfueco.2022.100078}{https://doi.org/10.1016/j.jfueco.2022.100078}

\bibitem{yun2021rowhammer}
Yun D, Park M, Bak G, Baeg S, Wen SJ. Exploitations of Multiple Rows Hammering and Retention Time Interactions in DRAM Using X-Ray Radiation. IEEE Access. 2021;9:137514-137523.
\href{https://doi.org/10.1109/access.2021.3117601}{https://doi.org/10.1109/access.2021.3117601}

\bibitem{kong2025twoT}
Kong S, Shim W. Design of 2T DRAM Cell With Surrounding Poly-Si Capacitor for Enhanced Retention and Mitigated Coupling Effect. IEEE Trans. Device Mater. Relib. 2025;25:460-464.
\href{https://doi.org/10.1109/tdmr.2025.3578692}{https://doi.org/10.1109/tdmr.2025.3578692}

\bibitem{kong2024pim1}
Hwan Kong S, Shim W. Advanced 2T0C DRAM Technologies for Processing-in-Memory-Part I: Vertical Transistor on Gate (VTG) DRAM Cell Structure. IEEE Trans. Electron Devices. 2024;71:6633-6638.
\href{https://doi.org/10.1109/ted.2024.3447612}{https://doi.org/10.1109/ted.2024.3447612}

\bibitem{kong2026threeT}
Kim K, Shim W. A novel 3T DRAM cell with depletion mode PMOS for extended retention time in processing-in-memory applications. Semicond. Sci. Technol. 2026;41:045011.
\href{https://doi.org/10.1088/1361-6641/ae544d}{https://doi.org/10.1088/1361-6641/ae544d}

\bibitem{nano2026gidl}
Ma H, Gu S, Zhai M, Liu H. Suppressing Gate-Induced Drain Leakage with an Asymmetric Gate Design in HiPco CNT FETs. Nanomaterials. 2026;16:653.
\href{https://doi.org/10.3390/nano16110653}{https://doi.org/10.3390/nano16110653}

\bibitem{tcad2026refresh}
Golman R, Cohen H, Shapiro G, Teman A. Cyclic Refresh Algorithm for Ensuring 100\% Memory Availability in Gain-Cell Embedded DRAM Macros. IEEE Trans. Comput.-Aided Des. Integr. Circuits Syst. 2026;45:4532-4543.
\href{https://doi.org/10.1109/tcad.2026.3652963}{https://doi.org/10.1109/tcad.2026.3652963}

\bibitem{ted2026refresh}
Ko K, Moon D, Kim J, Oh J, Park S, Song J. Dipole-Based Characteristic Enhancement Strategy for Refresh Performance Enhancement in DRAM BCAT Structure. IEEE Trans. Electron Devices. 2026;73:4885-4890.
\href{https://doi.org/10.1109/ted.2026.3708805}{https://doi.org/10.1109/ted.2026.3708805}

\bibitem{luza2022neutron}
Matana Luza L, Soderstrom D, Puchner H, Alia RG, Letiche M, Cazzaniga C, et al. Neutron-induced effects on a self-refresh DRAM. Microelectronics Reliability. 2022;128:114406.
\href{https://doi.org/10.1016/j.microrel.2021.114406}{https://doi.org/10.1016/j.microrel.2021.114406}

\bibitem{afacan2021review}
Afacan E, Lourenco N, Martins R, Dundar G. Review: Machine learning techniques in analog/RF integrated circuit design, synthesis, layout, and test. Integration. 2021;77:113-130.
\href{https://doi.org/10.1016/j.vlsi.2020.11.006}{https://doi.org/10.1016/j.vlsi.2020.11.006}

\bibitem{mina2022review}
Mina R, Jabbour C, Sakr GE. A Review of Machine Learning Techniques in Analog Integrated Circuit Design Automation. Electronics. 2022;11:435.
\href{https://doi.org/10.3390/electronics11030435}{https://doi.org/10.3390/electronics11030435}

\bibitem{lberni2024eadl}
Lberni A, Marktani MA, Ahaitouf A, Ahaitouf A. Analog circuit sizing based on Evolutionary Algorithms and deep learning. Expert Systems with Applications. 2024;237:121480.
\href{https://doi.org/10.1016/j.eswa.2023.121480}{https://doi.org/10.1016/j.eswa.2023.121480}

\bibitem{tang2024cnn}
Tang C, Chen X, Luo Y, Zeng Y. Automatic Optimal Design Method for Circuit Sizing Based on CNN Surrogate Model Assisted Differential Evolution Algorithm. IEEE Access. 2024;12:136238-136247.
\href{https://doi.org/10.1109/access.2024.3462952}{https://doi.org/10.1109/access.2024.3462952}

\bibitem{song2022mlsizing}
Song LY, Chou CY, Kuo TC, Liu CN, Huang JD. Machine Learning Assisted Circuit Sizing Approach for Low-Voltage Analog Circuits with Efficient Variation-Aware Optimization. ACM Trans. Des. Autom. Electron. Syst. 2022;28:1-22.
\href{https://doi.org/10.1145/3567422}{https://doi.org/10.1145/3567422}

\bibitem{ns2024pareto}
NS KS, Li P. Pareto Optimization of Analog Circuits Using Reinforcement Learning. ACM Trans. Des. Autom. Electron. Syst. 2024;29:1-14.
\href{https://doi.org/10.1145/3640463}{https://doi.org/10.1145/3640463}

\bibitem{vlsi2026transformer}
Sun M, Li H, Chen W, Li X, Li J, Luo Z, et al. Transformer-based surrogate modeling for multi-circuit analog sizing via Bayesian optimization. Integration. 2026;111:102809.
\href{https://doi.org/10.1016/j.vlsi.2026.102809}{https://doi.org/10.1016/j.vlsi.2026.102809}

\bibitem{tcad2026multitask}
Chen M, Hao Y, Gandara M, Imran M, Liu B. Trade-off-aware Analog Circuit Sizing Based on A Multitask Surrogate Model-assisted Evolutionary Algorithm. IEEE Trans. Comput.-Aided Des. Integr. Circuits Syst. 2026:1-1.
\href{https://doi.org/10.1109/tcad.2026.3677769}{https://doi.org/10.1109/tcad.2026.3677769}

\bibitem{akbar2022devsim}
Akbar C, Li Y, Thoti N. Device-Simulation-Based Machine Learning Technique for the Characteristic of Line Tunnel Field-Effect Transistors. IEEE Access. 2022;10:53098-53107.
\href{https://doi.org/10.1109/access.2022.3174685}{https://doi.org/10.1109/access.2022.3174685}

\bibitem{guo2023pinn}
Guo G, You H, Li C, Tang Z, Li O. A Physics-Informed Automatic Neural Network Generation Framework for Emerging Device Modeling. Micromachines. 2023;14:1150.
\href{https://doi.org/10.3390/mi14061150}{https://doi.org/10.3390/mi14061150}

\bibitem{liang2021esd}
Liang W, Yang X, Miao M, Loiseau A, Mitra S, Gauthier R. Novel ESD Compact Modeling Methodology Using Machine Learning Techniques for Snapback and Non-Snapback ESD Devices. IEEE Trans. Device Mater. Relib. 2021;21:455-464.
\href{https://doi.org/10.1109/tdmr.2021.3116599}{https://doi.org/10.1109/tdmr.2021.3116599}

\bibitem{yoon2023sram}
Yoon T, Jeong H. Machine Learning-Based Read Access Yield Estimation and Design Optimization for High-Density SRAM. IEEE Trans. Comput.-Aided Des. Integr. Circuits Syst. 2023;42:2618-2630.
\href{https://doi.org/10.1109/tcad.2022.3225066}{https://doi.org/10.1109/tcad.2022.3225066}

\bibitem{shaik2024sramaging}
Shaik JB, Guo X, Singhal S. Impact of Aging and Process Variability on SRAM-Based In-Memory Computing Architectures. IEEE Trans. Circuits Syst. I. 2024;71:2696-2708.
\href{https://doi.org/10.1109/tcsi.2024.3381935}{https://doi.org/10.1109/tcsi.2024.3381935}

\bibitem{touloupas2022mixedbo}
Touloupas K, Sotiriadis PP. Mixed-Variable Bayesian Optimization for Analog Circuit Sizing through Device Representation Learning. Electronics. 2022;11:3127.
\href{https://doi.org/10.3390/electronics11193127}{https://doi.org/10.3390/electronics11193127}

\bibitem{wang2022yieldbo}
Wang X, Yan C, Ma Y, Yu B, Yang F, Zhou D, et al. Analog Circuit Yield Optimization via Freeze-Thaw Bayesian Optimization Technique. IEEE Trans. Comput.-Aided Des. Integr. Circuits Syst. 2022;41:4887-4900.
\href{https://doi.org/10.1109/tcad.2022.3149723}{https://doi.org/10.1109/tcad.2022.3149723}

\bibitem{zhang2021nsga2mat}
Zhang P, Qian Y, Qian Q. Multi-objective optimization for materials design with improved NSGA-II. Materials Today Communications. 2021;28:102709.
\href{https://doi.org/10.1016/j.mtcomm.2021.102709}{https://doi.org/10.1016/j.mtcomm.2021.102709}

\bibitem{borisut2023alhs}
Borisut P, Nuchitprasittichai A. Adaptive Latin Hypercube Sampling for a Surrogate-Based Optimization with Artificial Neural Network. Processes. 2023;11:3232.
\href{https://doi.org/10.3390/pr11113232}{https://doi.org/10.3390/pr11113232}

\bibitem{phromphan2024lhs}
Phromphan P, Suvisuthikasame J, Kaewmongkol M, Chanpichitwanich W, Sleesongsom S. A New Latin Hypercube Sampling with Maximum Diversity Factor for Reliability-Based Design Optimization of HLM. Symmetry. 2024;16:901.
\href{https://doi.org/10.3390/sym16070901}{https://doi.org/10.3390/sym16070901}

\bibitem{liu2025xaimat}
Liu B, Liu P, Lu W, Olofsson T. Explainable Artificial Intelligence (XAI) for Material Design and Engineering Applications: A Quantitative Computational Framework. Int Journal of Mech Sys Dyn. 2025;5:236-265.
\href{https://doi.org/10.1002/msd2.70017}{https://doi.org/10.1002/msd2.70017}

\bibitem{lundberg2017shap}
Lundberg SM, Lee SI. A unified approach to interpreting model predictions. Advances in Neural Information Processing Systems (NeurIPS). 2017;30:4765-4774.

\bibitem{lundberg2020treeshap}
Lundberg SM, Erion G, Chen H, DeGrave A, Prutkin JM, Nair B, et al. From local explanations to global understanding with explainable AI for trees. Nat Mach Intell. 2020;2:56-67.
\href{https://doi.org/10.1038/s42256-019-0138-9}{https://doi.org/10.1038/s42256-019-0138-9}

\bibitem{apley2020ale}
Apley DW, Zhu J. Visualizing the Effects of Predictor Variables in Black Box Supervised Learning Models. Journal of the Royal Statistical Society Series B: Statistical Methodology. 2020;82:1059-1086.
\href{https://doi.org/10.1111/rssb.12377}{https://doi.org/10.1111/rssb.12377}

\bibitem{ruhe2023dtwin}
Ruhe S, Schaefer K, Branz S, Nicolai S, Bretschneider P, Westermann D. Design and Implementation of a Hierarchical Digital Twin for Power Systems Using Real-Time Simulation. Electronics. 2023;12:2747.
\href{https://doi.org/10.3390/electronics12122747}{https://doi.org/10.3390/electronics12122747}

\bibitem{massaro2025dtwin}
Massaro A. Electronic Artificial Intelligence-Digital Twin Model for Optimizing Electroencephalogram Signal Detection. Electronics. 2025;14:1122.
\href{https://doi.org/10.3390/electronics14061122}{https://doi.org/10.3390/electronics14061122}

\bibitem{mckay1979lhs}
McKay MD, Beckman RJ, Conover WJ. A Comparison of Three Methods for Selecting Values of Input Variables in the Analysis of Output from a Computer Code. Technometrics. 1979;21:239.
\href{https://doi.org/10.2307/1268522}{https://doi.org/10.2307/1268522}

\bibitem{zhao2006ptm}
Zhao W, Cao Y. New Generation of Predictive Technology Model for Sub-45 nm Early Design Exploration. IEEE Trans. Electron Devices. 2006;53:2816-2823.
\href{https://doi.org/10.1109/ted.2006.884077}{https://doi.org/10.1109/ted.2006.884077}

\bibitem{pelgrom1989matching}
Pelgrom MJM, Duinmaijer ACJ, Welbers APG. Matching properties of MOS transistors. IEEE J. Solid-State Circuits. 1989;24:1433-1439.
\href{https://doi.org/10.1109/jssc.1989.572629}{https://doi.org/10.1109/jssc.1989.572629}

\bibitem{ngspice}
Vogt H, Nenzi P, Warning D, et al. Ngspice, the open-source mixed-level / mixed-signal circuit simulator (version 46). 2024. \href{https://ngspice.sourceforge.io}{https://ngspice.sourceforge.io}

\bibitem{saltelli2010sobol}
Saltelli A, Annoni P, Azzini I, Campolongo F, Ratto M, Tarantola S. Variance based sensitivity analysis of model output. Design and estimator for the total sensitivity index. Computer Physics Communications. 2010;181:259-270.
\href{https://doi.org/10.1016/j.cpc.2009.09.018}{https://doi.org/10.1016/j.cpc.2009.09.018}

\bibitem{chen2016xgboost}
Chen T, Guestrin C. XGBoost: a scalable tree boosting system. Proceedings of the 22nd ACM SIGKDD International Conference on Knowledge Discovery and Data Mining. 2016:785-794.
\href{https://doi.org/10.1145/2939672.2939785}{https://doi.org/10.1145/2939672.2939785}

\bibitem{ke2017lightgbm}
Ke G, Meng Q, Finley T, Wang T, Chen W, Ma W, et al. LightGBM: a highly efficient gradient boosting decision tree. Advances in Neural Information Processing Systems (NeurIPS). 2017;30:3146-3154.

\bibitem{prokhorenkova2018catboost}
Prokhorenkova L, Gusev G, Vorobev A, Dorogush AV, Gulin A. CatBoost: unbiased boosting with categorical features. Advances in Neural Information Processing Systems (NeurIPS). 2018;31:6638-6648.

\bibitem{pedregosa2011sklearn}
Pedregosa F, Varoquaux G, Gramfort A, Michel V, Thirion B, Grisel O, et al. Scikit-learn: machine learning in Python. Journal of Machine Learning Research. 2011;12:2825-2830.

\bibitem{deb2002nsga2}
Deb K, Pratap A, Agarwal S, Meyarivan T. A fast and elitist multiobjective genetic algorithm: NSGA-II. IEEE Trans. Evol. Computat. 2002;6:182-197.
\href{https://doi.org/10.1109/4235.996017}{https://doi.org/10.1109/4235.996017}

\bibitem{deb2014nsga3}
Deb K, Jain H. An Evolutionary Many-Objective Optimization Algorithm Using Reference-Point-Based Nondominated Sorting Approach, Part I: Solving Problems With Box Constraints. IEEE Trans. Evol. Computat. 2014;18:577-601.
\href{https://doi.org/10.1109/tevc.2013.2281535}{https://doi.org/10.1109/tevc.2013.2281535}

\bibitem{zhang2007moead}
Qingfu Zhang, Hui Li. MOEA/D: A Multiobjective Evolutionary Algorithm Based on Decomposition. IEEE Trans. Evol. Computat. 2007;11:712-731.
\href{https://doi.org/10.1109/tevc.2007.892759}{https://doi.org/10.1109/tevc.2007.892759}

\bibitem{storn1997de}
Storn R, Price K. Differential Evolution - A Simple and Efficient Heuristic for global Optimization over Continuous Spaces. Journal of Global Optimization. 1997;11:341-359.
\href{https://doi.org/10.1023/a:1008202821328}{https://doi.org/10.1023/a:1008202821328}

\bibitem{kennedy1995pso}
Kennedy J, Eberhart R. Particle swarm optimization. Proceedings of ICNN'95 - International Conference on Neural Networks. 1995;4:1942-1948.
\href{https://doi.org/10.1109/icnn.1995.488968}{https://doi.org/10.1109/icnn.1995.488968}

\bibitem{bergstra2011tpe}
Bergstra J, Bardenet R, Bengio Y, Kegl B. Algorithms for hyper-parameter optimization. Advances in Neural Information Processing Systems (NeurIPS). 2011;24:2546-2554.

\bibitem{akiba2019optuna}
Akiba T, Sano S, Yanase T, Ohta T, Koyama M. Optuna: a next-generation hyperparameter optimization framework. Proceedings of the 25th ACM SIGKDD International Conference on Knowledge Discovery and Data Mining. 2019:2623-2631.
\href{https://doi.org/10.1145/3292500.3330701}{https://doi.org/10.1145/3292500.3330701}

\bibitem{blank2020pymoo}
Blank J, Deb K. Pymoo: Multi-Objective Optimization in Python. IEEE Access. 2020;8:89497-89509.
\href{https://doi.org/10.1109/access.2020.2990567}{https://doi.org/10.1109/access.2020.2990567}
